%==============================================================================
\chapter{Enveloppe budgétaire et ressources}
%==============================================================================

\section{Cadre}
Le projet s'inscrit dans un cadre \textbf{académique (mini-projet)}. L'enveloppe
budgétaire se traduit donc principalement en \textbf{ressources humaines,
logicielles et matérielles} mobilisées, sans coût financier direct.

\section{Ressources mobilisées}
\begin{xltabular}{\textwidth}{>{\bfseries}p{4cm} X}
\toprule
\textbf{Ressource} & \textbf{Détail} \\
\midrule
\endhead
Humaines   & Équipe projet étudiante (conception, développement, tests, documentation). \\
Logicielles & Environnement d'exécution conteneurisé, SGBD relationnel, IDE, outils UML, chaîne de compilation. \\
Matérielles & Postes de développement, capteurs simulés pour le volet indicateurs. \\
Externes   & Compte OAuth Google, certificats X.509 (auto-signés en démonstration). \\
\bottomrule
\end{xltabular}

\begin{encadre}
\textbf{Contrainte académique majeure :} tout usage de générateurs de code
(ChatGPT, DeepSeek et dérivés) est proscrit et sera considéré comme une fraude
à l'épreuve d'examen.
\end{encadre}

%------------------------------------------------------------------------------
\section{Valorisation économique : hypothèse d'industrialisation}
%------------------------------------------------------------------------------
La section précédente décrit le coût \emph{réel} du projet, nul en dépenses
directes. Par exercice de valorisation, cette section estime ce que coûterait le
même produit \textbf{s'il était réalisé par une société de services}. Les
montants sont des \textbf{hypothèses de travail}, non des dépenses engagées.

\subsection{Charge de référence}
Le plan de charge du chapitre des délais retient \textbf{304 story points} sur
8~sprints. À $0{,}79$~j-h par point, la charge s'établit à $\approx 240$~j-h ;
à $1$~j-h par point --- hypothèse prudente pour une équipe qui découvre la pile
--- elle atteint $\approx 304$~j-h. C'est cette fourchette qui sert de base au
chiffrage.

\subsection{Répartition d'une enveloppe de référence}
Pour une enveloppe de \textbf{200\,000~USD}, la ventilation suivante reflète les
usages du marché. La part de développement n'est pas majoritaire : c'est une
erreur classique de chiffrage que de l'assimiler au coût total.

\begin{xltabular}{\textwidth}{>{\bfseries}p{5.2cm} r r X}
\toprule
\textbf{Poste} & \textbf{\%} & \textbf{USD} & \textbf{Justification} \\
\midrule
\endhead
Développement & 45 & 90\,000 & Back-end, front-end, intégration \\
Gestion de projet / analyse & 12 & 24\,000 & Poste le plus souvent sous-estimé \\
Qualité et tests & 10 & 20\,000 & Couverture $\geq 70\,\%$, tests d'intégration et de charge \\
Design UX/UI et accessibilité & 8 & 16\,000 & Coût réduit : la charte existe déjà \\
Infrastructure et outillage (an 1) & 6 & 12\,000 & Hébergement, sauvegardes, supervision \\
Audit de sécurité externe & 5 & 10\,000 & Données de géolocalisation sensibles \\
Provision pour aléas & 14 & 28\,000 & Sans réserve, une enveloppe de cette taille dérape \\
\bottomrule
\end{xltabular}

\subsection{Sensibilité au taux journalier}
Le facteur dominant n'est pas technique mais géographique. À un taux moyen de
600--900~USD par jour-homme, l'enveloppe finance $\approx 220$--$280$~j-h :
elle couvre tout juste la charge de référence, \textbf{sans marge}. À un taux de
250--400~USD, elle finance $\approx 500$--$700$~j-h et autorise une équipe de
cinq à six personnes ainsi qu'une provision saine.

\subsection{Périmètre ajustable}
Si l'enveloppe se tend, l'ordre de retrait suivant préserve le c\oe ur du
produit : d'abord les fonctions avancées (annexe~: suivi de reine, QR code de
lot), puis la détection d'anomalie adaptative --- dont le calibrage sans
historique réel constitue un risque identifié ---, enfin l'ingestion réelle des
capteurs, le modèle de données suffisant en première version. Restent
incompressibles les entités métier, les visites, l'authentification et le
contrôle d'accès, le mode hors-ligne et les tableaux de bord.

\subsection{Coût total de possession}
Le coût de construction ne représente qu'une part du coût de vie du produit. La
maintenance corrective et évolutive se chiffre usuellement à
\textbf{15--20\,\% par an} de l'enveloppe de construction, soit
30\,000--40\,000~USD annuels. Sur trois ans, le coût total de possession
approche donc \textbf{300\,000~USD}. S'y ajoutent les postes régulièrement omis
d'un chiffrage initial : reprise des données existantes de l'exploitant,
formation et conduite du changement auprès d'utilisateurs de terrain,
traduction professionnelle relue pour les trois langues, et mise en conformité
relative aux données à caractère personnel.

\begin{encadre}
\textbf{Portée de cette section :} il s'agit d'une \emph{valorisation}
pédagogique destinée à situer l'effort du projet dans les ordres de grandeur du
marché. Elle n'engage aucune dépense et ne modifie pas le cadre académique
défini en début de chapitre.
\end{encadre}

